% Exact Vardi 1982 proof/sketch boundaries and explicit reconstruction.

\begin{annexitem}{annex:p9.n-c-004.2.theorem5}{thm:p09-vardi82-theorem5}{Two uniform reductions}
\begin{proof}
The membership bounds are the positive-iteration lemma below. We supply
both lower-bound reductions, including their logspace uniformity. Vardi's
printed proof invokes Path System Accessibility for the first and sketches
a Boolean-addressed tableau for the second
\cite[Theorem~5, pp.~141--142]{Vardi1982}. The details here are a handbook
reconstruction, not additional text attributed to the extended abstract.

\paragraph{The fixed query and its derivations.}
A path system consists of a finite set of points, a unary set \(A\) of
initial points, a ternary relation \(P\), and a target \(g\). Its accessible
points are the least set containing \(A\) and closed under
\(x,y\mapsto z\) whenever \(P(x,y,z)\). One fixed expression is
\[
 \Upsilon X.z.\bigl(A(z)\lor
      \exists x\exists y(X(x)\land X(y)\land P(x,y,z))\bigr).
\]
Induction on stages identifies its members with points having finite
derivation trees: initial points are the leaves, and an internal node has
the two premises of one triple. Conversely, induction on a derivation
tree's height places its root in a finite stage. Iteration adds at most
one point per strict increase, so accessibility is computable in polynomial
time. It remains to establish hardness without assuming it from the name
of the problem.

\paragraph{From a polynomial-time machine to one path system.}
Fix any language in \(PTIME\) and a deterministic machine deciding it.
A standard one-tape simulation suffices with polynomial overhead: represent
each of a fixed number of tapes and its head mark on tracks of one tape,
scan the visited interval to collect the scanned symbols, and scan it again
to perform the finitely many writes and head moves. After \(t\) simulated
steps the interval has length \(O(t+n+1)\), so the overhead is polynomial.
Thus fix a one-tape machine with a polynomial bound \(T=T(n)\geq n+2\),
and extend both halting states by stationary, unchanged transitions.
Use \(T+2\) squares; enlarge \(T\) if necessary so that a genuine run
does not encounter the right boundary. A left-boundary move follows the
machine's specified convention.

Let \(\Omega\) be the fixed finite set of complete cell tokens: either
a tape symbol without a head, or a tape symbol with the head and its
control state. A genuine next-row token at a square depends only on the
three adjacent old tokens. There is consequently a fixed map
\(f:\Omega^3\to\Omega\), with separate endpoint cases using a virtual
outside token. To verify this, a head in the central cell writes there
and either stays or moves; a head in a neighbour can enter the cell with
its new state; in the absence of these events the symbol is unchanged.
Assign arbitrary values to inconsistent triples with competing heads;
the induction below never uses them in a genuine row.

Create a point \(v_{t,p,a}\) for each time \(0\leq t\leq T\), square
\(p\), and token \(a\). Put precisely the genuine initial-row points
in \(A\). For an interior square and each token triple
\((a_-,a,a_+)\), create an additional point \(u_{t,p,a_-,a}\) and put
the following two triples into \(P\):
\[
\begin{split}
 &(v_{t,p-1,a_-},v_{t,p,a},u_{t,p,a_-,a}),\\
 &(u_{t,p,a_-,a},v_{t,p+1,a_+},
                  v_{t+1,p,f(a_-,a,a_+)}).
\end{split}
\]
Only times \(t<T\) occur here. At an endpoint the local update uses two
available cell tokens; put their pair directly into a triple with the
new token. If an update has only one premise, repeat it in both premise
slots. Finally, for every accepting-head token at any time and square,
put \((v_{t,p,a},v_{t,p,a},g)\) into \(P\), with \(g\) a fresh point.

Induction on time shows that exactly the genuine token point at each
square becomes accessible. The intermediate point requires its two actual
old tokens, and the second triple requires the third one; determinism
then forces the correct new token. Conversely those three old points
derive that token. Endpoints satisfy the same argument with the virtual
token fixed. No other triples produce token points or \(g\). Thus \(g\)
is accessible exactly when the machine accepts the input word.

This is a logspace many-one reduction. Give the points consecutive
identifiers by their fixed-shape index tuples and enumerate all the
displayed triples with counters for \(t,p\) and the finitely many tokens.
The counters and each identifier have \(O(\log(n+2))\) bits because
\(T\) is a fixed polynomial. Initialization is found by rescanning the
input at the required position; the reduction never simulates the run.
There are polynomially many points and triples. The output is the
explicit database \((D,A,P)\) together with target \(g\), and the query
expression above is independent of the input and of its length. This
proves D-hardness and also supplies the Path System Accessibility
completeness ingredient cited, but not proved, in the source.

\paragraph{The fixed Boolean database and variable-width addresses.}
For E-hardness fix a language in \(EXPTIME\) and a deterministic decider.
The same tape conversion has at most exponential time after enlarging a
fixed polynomial exponent. Modify acceptance to return the head to square
zero, write a designated symbol there, and enter a stationary accepting
state; rejection has a different stationary state. The return takes at
most the number of previously visited squares. Hence there is a fixed
integer polynomial \(q\) such that, on an input \(w\) of length \(n\),
all this takes fewer than \(2^{q(n)}\) steps and visits fewer than that
many squares. Set \(N=q(n)+2\), increasing the constant term so that
\(N\geq1\) and \(2^N>n+2\).

Keep \(B_0=(\{0,1\},O)\), \(O=\{1\}\), fixed. A Boolean tuple of
length \(N\) codes an integer in \([0,2^N-1]\), most significant bit
first. Define \(Z_N(\bar x)\) by \(\bigwedge_i\neg O(x_i)\), and
\(U_N(\bar x)\) by \(\bigwedge_i O(x_i)\). The formula
\[
 S_N(\bar x,\bar y)=
 \bigvee_{i=1}^{N}\left(
    \bigwedge_{j<i}x_j=y_j\ \land\ \neg O(x_i)\land O(y_i)
    \ \land\ \bigwedge_{j>i}(O(x_j)\land\neg O(y_j))\right)
\]
is ordinary nonwrapping binary successor. Its length is \(O(N^2)\),
even with no abbreviations shared between subformulas. Tests for any
specified address are conjunctions of \(N\) literals of \(O\).
Choose a fixed width \(m\) coding the finite token alphabet \(\Omega\)
injectively in \(\{0,1\}^m\). Write \(C_a(\bar z)\) for the
corresponding fixed conjunction of bit tests.

The only recursively defined relation is
\(R(\bar t,\bar p,\bar z)\), of rank \(2N+m\). Its initial condition
requires \(Z_N(\bar t)\). At each position \(i<n\), an address test
selects the token for the literal input letter \(w_i\), with the initial
head-state token at position zero. Positions outside that finite list
receive blank tokens; if \(n=0\), position zero receives the initial
head on a blank square. These conditions are a first-order formula of
length \(O((n+1)N)\), with no occurrence of \(R\).

For each fixed interior token triple, the successor disjunct has free
arguments \((\bar u,\bar p,\bar z)\) and existentially quantified
predecessor coordinates. It conjoins
\[
\begin{gathered}
 S_N(\bar t,\bar u),\quad
 S_N(\bar p_-,\bar p),\quad S_N(\bar p,\bar p_+),\quad
 C_{f(a_-,a,a_+)}(\bar z),\\
 R(\bar t,\bar p_-,\bar z_-),\quad C_{a_-}(\bar z_-),\quad
 R(\bar t,\bar p,\bar z_0),\quad C_a(\bar z_0),\quad
 R(\bar t,\bar p_+,\bar z_+),\quad C_{a_+}(\bar z_+).
\end{gathered}
\]
At the left or right endpoint use \(Z_N(\bar p)\) or
\(U_N(\bar p)\), omit the unavailable neighbour fact, and use the
fixed virtual-boundary token in \(f\). Disjoin these finitely many cases
and the initial condition to obtain \(\varphi_w(R,\bar u,\bar p,\bar z)\).
Every occurrence of \(R\) is positive; all remaining tests are
first-order over the fixed predicate \(O\). The expression
\[
 e_w=\Upsilon R.(\bar u,\bar p,\bar z).
                         \varphi_w(R,\bar u,\bar p,\bar z)
\]
therefore belongs to the original outer-operator grammar \(YF\).

An induction on the coded time proves soundness: every derived token is
the actual token in the deterministic run, since its predecessor facts
are genuine and the local update is deterministic. No predecessor of
time zero exists. The reverse induction on time proves completeness:
row zero occurs at stage one, and all tokens of the next row occur after
the preceding row. Thus the least fixed point is exactly the full
computation table through time \(2^N-1\). The extra clock and tape slack
ensures that boundary rejection cannot affect a genuine computation.
Because halting persists, the tuple
\[
 \bar c_w=(1^N,0^N,\operatorname{code}(a_{\rm acc}))
\]
belongs to \(Q_{e_w}(B_0)\) exactly when \(w\) is accepted. The
designated token \(a_{\rm acc}\) is the fixed accepting head and symbol
at square zero. No existential projection or \(YF^+\) extension is
needed to test this one output tuple.

Finally \(w\mapsto(\bar c_w,e_w)\) is logspace computable, not merely
polynomial-size. A transducer computes the fixed polynomial \(N\) with
\(O(\log(n+2))\)-bit counters. It emits each successor conjunction by
nested counters over bit positions, each input-address test by a counter
over \(i<n\) and a rescan for \(w_i\), and each local-update case from
the fixed machine transition table. Binary digits of an index and the
names of bound variables can be emitted within the same space; distinct
variable blocks avoid capture. Fully expanded output length is polynomial
in \(n\), since the number of token cases is fixed and the address
formulas have polynomial length. This proves E-hardness on the single
database \(B_0\). Together with the membership bounds, the two
completeness assertions follow.
\end{proof}
\end{annexitem}

\begin{annexitem}{annex:p9.n-c-004.2.iteration}{thm:p09-vardi82-iteration}{Iteration and input size}
\begin{proof}
Expand implications into Boolean connectives before checking polarity.
For \(X\subseteq X'\), use mutual formula induction: a formula positive
in the relation variable preserves truth from \(X\) to \(X'\), while a
negative formula preserves truth from \(X'\) to \(X\).
The relation atom gives the positive inclusion; other atoms are unchanged.
Conjunction and disjunction preserve each induction direction, and the
quantifiers do so by retaining all witnesses or all tested domain values.
Negation interchanges the two induction directions and their polarities.
This covers every formula constructor.  Thus the operator
\(\Phi_B\) is monotone.

Put \(X_0=\varnothing\), \(X_{j+1}=\Phi_B(X_j)\).
Then \(X_0\subseteq X_1\), and induction using monotonicity gives
\(X_j\subseteq X_{j+1}\) at every stage.  Each strict increase adds a new
tuple among the \(\lvert D\rvert^r\) possible tuples.  After at most that
many increases there is equality at two consecutive stages, and the equality
makes the relation a fixed point.  If \(Y\) is any fixed point, induction
gives \(X_j\subseteq Y\), starting with the empty relation.  The stable
relation is consequently the least fixed point.  This proves existence,
termination and the claimed bound, including an empty domain at positive rank.

For a fixed matrix with \(w\) individual-variable slots, evaluate it by
enumerating assignments and their prefixes, at most
\((1+\lvert D\rvert)^{w+1}\), respecting
universal conjunction and existential disjunction.  Store the current
relation table and scan an explicitly stored database when necessary.
Each table lookup costs at most its polynomial-size table scan.
Both \(r\) and \(w\) are fixed; at most \(\lvert D\rvert^r+1\)
iterations therefore take polynomial time in the explicit input size.
The candidate output tuple is checked against the final table.

If expression size \(s\) varies, its arities and number of variable slots
are at most \(s\) under explicit expression encoding.  With a fixed
database of \(d\) elements, the table and assignment counts are bounded by
\((1+d)^{O(s)}\), up to polynomial expression-processing overhead.
For \(d=0\) or \(1\) the bound only improves; for fixed \(d\geq2\)
it is exponential in \(s\).  For combined explicit input size \(N\),
both \(d\) and \(s\) are at most \(N\), so the same simple bound is
at most \(2^{O(N\log N)}\), with a larger polynomial exponent if the
chosen elementary implementation requires it.  First-order expressions
without the outer operator have the corresponding cheaper assignment bound.
Equality restrictions and a fixed number of projections preserve polynomial
data complexity.  None of these calculations turns a formula-dependent rank
into a fixed constant on the expression or combined axes.
\end{proof}
\end{annexitem}

\begin{annexitem}{annex:p9.n-c-004.2.ordered}{thm:p09-vardi82-remark2}{Deterministic and existential tableaux at their exact query scope}
\begin{proof}
The source's Remark~2 states the two capture mechanisms briefly. We
reconstruct the deterministic construction within its single-operator
grammar, then derive the full existential-second-order query statement
from the proved Fagin theorem. The empty-universe cases are treated
directly, without an unqualified appeal to a normal-form theorem.

\paragraph{What order and determinism replace.}
The source first invokes the existential-second-order route: existential
relations choose a domain order and an accepting nondeterministic
computation.  With the order supplied by the database and the machine
deterministic, neither guess is needed; positive dependencies can construct
the unique computation.  The subsequent equality restrictions and existential
projection explain the extra \(+\) in \(YF^+\)
\cite[Remark~2, pp.~142--143]{Vardi1982}.
The argument below uses a single defined relation with a first-order positive
matrix, not an assumed closure under arbitrary nested least fixed points.

\paragraph{First-order addresses, tags and a fixed query.}
Fix a total consistent polynomial-time query of output rank \(b\).
For a partial query, this same construction decides its recognition graph
and thus produces its total graph-equivalent completion, not its original
domain of definition.
Let the database's input order be a weak total order \(\leq\), and
define \(x<y\) by \(x\leq y\land x\ne y\).
Minimum, maximum and immediate successor have first-order definitions:
\[
\begin{aligned}
 \operatorname{Min}(x)&:=\forall y(x\leq y),&
 \operatorname{Max}(x)&:=\forall y(y\leq x),\\
 \operatorname{Succ}(x,y)&:=x<y\land\neg\exists z(x<z\land z<y).
\end{aligned}
\]
On \(D^k\), lexicographic order is the finite disjunction over a first
unequal coordinate, with all preceding coordinates equal; immediate tuple
successor uses this order and excludes an intermediate tuple.  Unlike the
source's earlier guessed-order argument, these formulas use the \emph{given}
\(\leq\).  No tuple address wraps around.

For \(\lvert D\rvert\geq2\), minimum and maximum are distinct and
can represent Boolean 0 and 1.  Any fixed finite set of machine symbols has
an injective code into Boolean tuples of a fixed width \(w\): take distinct
words in \(\{0,1\}^w\), and test each word with equalities to the
first-order-defined endpoints.  This avoids falsely requiring that a
two-element universe itself contain every state or tape-symbol code.
Each code test is a fixed first-order formula; no endpoint constant has to be
in the input vocabulary.

All query tuples are simulated at once.  A prefix \(\bar d\in D^b\)
of the defined relation separates their tableaux.  Every predecessor fact
used to make a new fact has exactly this same prefix.  Thus these slots are
part of the operator's bound output tuple, as required by the source's
definition, rather than undeclared free parameters.

\paragraph{An explicitly definable machine input.}
Consistency permits replacing the finite ordered domain by its order ranks
\(0,\ldots,d-1\).  Relabeling all relations and the candidate tuple this
way preserves query membership.  Under the standard explicit database
encoding this representative has length polynomial in \(d=\lvert D\rvert\)
for fixed schema and output rank.  A deterministic query-graph recognizer
therefore has a polynomial-in-\(d\) bound on these representatives even
when the original numerical names were very long.  This step does not
assume that a query may inspect arbitrary numerical names in violation of
the consistency criterion.

For initialization one can use a padded input word whose symbols are
first-order definable directly in the ordered database.  Here is a concrete
layout.  Let \(h\) dominate all input relation ranks and let a fixed
Boolean tag tuple select blocks.  An address consists of that tag and
\(h\) further slots.  Reserve one block for a unary count of the domain:
put a universe marker at addresses whose first \(h-1\) slots are minimum,
letting the final slot vary.  This gives exactly \(d\) markers.
For each input relation \(R_i\), reserve a block and put its 0/1 truth value
at addresses whose first \(h-a_i\) slots are minimum and whose last
\(a_i\) slots are its arguments.  For each candidate coordinate \(d_j\),
reserve a block with a marker exactly where the last slot equals \(d_j\)
and the other argument slots are minimum.  Fill all unused addresses with
a separator symbol and place an end marker at the last full address, after
choosing sufficiently many tag slots that this address is an unused block.
Every symbol test in this layout is first-order.  Its total length is
\(d^c\) for a fixed \(c\); it includes the exact universe, relation tables
and candidate ranks, without assuming a built-in binary-bit predicate.

A fixed deterministic machine can scan this word, count the universe markers,
extract the relation bits and candidate-marker positions, and form the
standard ordered-rank encoding before running the query recognizer.  All
counts and work intervals are polynomial in \(d\).  A usual finite-track
one-tape simulation scans the visited work interval to find head marks and
then implements each fixed transition; its overhead is polynomial.  We may
therefore use a deterministic one-tape machine with a polynomial-in-\(d\)
time bound.  The elementary finite-track transformation can also eliminate
stationary moves by two moves with additional control states.  For the tableau
below stationary moves could instead be included in the same local map.

Choose a fixed \(k\) large enough that, for every \(d\geq2\), this
entire input and computation fit into \(d^k\) tape squares and
\(d^k-1\) transitions.  If a preliminary bound is \(C d^a\), increasing
\(k\) by \(\lceil\log_2(C+1)\rceil+2\) supplies this slack for all
\(d\geq2\).  Rejection and acceptance states persist to the clock's end.
The ordered tuple addresses serve as both times and tape squares.

\paragraph{Positive facts, including negative cell information.}
Let \(\Omega\) be the finite alphabet of cell tokens: ordinary tape symbols
and tape symbols paired with a control state to mark the unique head.
The whole token states a complete cell value.  A token saying that the cell
contains 0, or contains a blank without a head, is still a \emph{positive}
relation fact.  This removes any need to negate the evolving relation in
order to discover absent tape bits or absent heads.

Use one relation
\(T(\bar d,\bar t,\bar p,\bar a,f_0,f_1)\), with \(k\)-tuple time
\(\bar t\), \(k\)-tuple cell position \(\bar p\), \(w\)-tuple
token code \(\bar a\), and two flag slots.  Computational facts have
\(f_0=\min D\), \(f_1=\max D\), hence distinct flags; output-signal
facts have equal flags.  The distinction is a first-order condition on
the arguments, not a negative test of \(T\).
Denote computational membership with a token \(a\in\Omega\) by
\(T_a(\bar d,\bar t,\bar p)\); this abbreviates a first-order existential
formula selecting the token code and the distinct endpoint flags, then
testing \(T\).  It is positive in \(T\).

The initialization clause, only for \(d\geq2\), places one computational
token at every square at the first time.  Embed the \(c\)-tuple input addresses above in \(k\)-tuple squares by
fixing the first \(k-c\) coordinates to the minimum. They occupy an
initial tuple interval; later squares are blank.  The head/state component
is present only at the first square.  These are fixed first-order cases.

For a deterministic one-tape machine there is a finite local update map
\(f:\Omega^3\to\Omega\).  A cell with no neighbouring head retains its
symbol; the head's old square receives the written symbol; its new square
receives the moved head and new control state.  At the two endpoints use
virtual outside-boundary tokens and the machine's prescribed rejection
behaviour, never a cyclic neighbour.  Define \(f\) arbitrarily on token
triples incompatible with a genuine one-head row.  Such triples never
occur in the derived computation.  For each actual token triple, the new
token clause is
\[
\begin{gathered}
 \operatorname{Succ}_k(\bar t,\bar u)\land
 T_{a_-}(\bar d,\bar t,\bar p_-)\land
 T_{a}(\bar d,\bar t,\bar p)\land
 T_{a_+}(\bar d,\bar t,\bar p_+)\\
 \longrightarrow\quad
 T_{f(a_-,a,a_+)}(\bar d,\bar u,\bar p).
\end{gathered}
\]
Here \(\bar p_-\), \(\bar p_+\) are tuple predecessor and successor;
the boundary cases omit the unavailable predecessor fact and substitute
the fixed virtual token.  To obtain a first-order operator matrix, disjoin
over the finitely many token triples and existentially quantify the
predecessor coordinates.  The displayed implication describes the clause;
the operator itself has the successor fact's argument variables free and
uses the conjunction on the left as its positive membership condition.

Union the initial facts with these successor conditions and the output-signal
conditions below.  This is one first-order formula positive in \(T\),
whose free variables are precisely the full bound tuple of \(T\).
No relation guessing, negation of \(T\), nested operator or unbounded
syntactic iteration is used.

\paragraph{Soundness and completeness of the derived rows.}
At the first operator stage every candidate prefix has exactly its genuine
initial row.  An induction on time says that any derived computational token
at that time is its genuine token: its predecessor facts are genuine and
the deterministic local map gives one token.  The initialization and
successor cases have different time conditions, and no predecessor of the
first time exists.  Output-signal facts cannot be used as computational
predecessors because their flags are equal.  Thus there are no spurious
computational tokens.

Conversely the same induction says every token of a genuine next row is
derived once all tokens of the preceding row have been derived.  It uses at
most three predecessor facts, or the stated virtual-boundary case.  All
cells of row~0 are obtained at stage~1; row~j is obtained by stage~j+1.
The chosen slack ensures that every accepting run has an accepting row
within the tuple time bound.  Therefore the fixed point contains an
accepting-head token for prefix \(\bar d\) exactly when that candidate
tuple belongs to the fixed query answer.

\paragraph{Answer projection and singleton inputs.}
The operator emits equal-flag output signals for prefix \(\bar d\) when
some computational fact for that prefix is an accepting-head token.  Its
other coordinates can be set equal to the minimum, with equal flags there,
by first-order tests inside the operator.  No signal is emitted otherwise.
For a singleton database there are finitely many input relation diagrams
at a fixed schema and only one candidate tuple of each fixed rank.  The
query machine decides those finitely many cases.  A first-order disjunction
of the positive and negative relation atoms of the accepted singleton
diagrams directly emits equal-flag signals for those cases.  The
computational initialization is guarded by \(\lvert D\rvert\geq2\),
so singleton code collisions cannot produce spurious computation.  Rejecting
singleton diagrams generate no facts.

The final query is the existential projection of the fixed-point expression
with the single equality restriction \(f_0=f_1\).  All the auxiliary
coordinates are projected away; the prefix \(\bar d\) remains.  This is
exactly the restriction/projection operation permitted by the source's
\(YF^+\) remark.  In particular projecting an accepting state is not
silently credited to the bare \(YF\) expression returning its entire
tableau relation.  For nonempty databases this proves the construction for
every fixed answer rank, including Boolean extraction.

For an empty domain and positive answer rank, the answer must be empty and
the construction also returns no tuples.  For rank zero a total query may
be true on the empty domain.  Let \(\epsilon\) be that one truth value,
decided by its fixed machine.  In the closed language \(Y\) the final
Boolean formula is
\[
 \bigl((\neg\exists x\,x=x)\land\epsilon\bigr)
 \lor \exists\bar t\exists\bar p\exists\bar a\exists f_0\exists f_1
   \bigl(f_0=f_1\land
      [\Upsilon T.\psi](\bar t,\bar p,\bar a,f_0,f_1)\bigr),
\]
where \(\epsilon\) is a first-order truth or falsity sentence and
\(\psi\) is the positive operator just constructed; the bracketed notation
abbreviates membership in the full-bound-tuple source expression
\(\Upsilon T.(\bar t,\bar p,\bar a,f_0,f_1).\psi\), not a new operator grammar.
On empty input the
existential tableau disjunct is false and the first disjunct is exactly
the intended answer; on nonempty input the first disjunct is false and
the projection is correct.  This is a \(Y\) expression by its stated Boolean and first-order
closure. If \(\epsilon\) is false, the original \(YF^+\) projection
already gives the right answer on every input. If \(\epsilon\) is true,
the additional empty disjunct is kept in \(Y\); it is not claimed to
have the same literal \(YF^+\) normal form. The separate
\cref{thm:p09-vardi82-empty-boundary} proves why that distinction is
necessary. The source's p.~145 normal-form assertion is therefore not
being extended to an unverified empty-domain convention.

For the converse, apply the positive-iteration lemma.
Each fixed \(YF^+\) expression has fixed relation rank and a fixed
number of projected coordinates. Closing fixed \(YF\) expressions
under finitely many Boolean connectives and first-order quantifiers still
takes polynomial data time: compute each fixed-point relation table, then
enumerate the finitely many fixed-length assignment tuples for the outer
formula. Thus the claimed bound also holds for \(Y\). The constructions
and evaluations commute with order-preserving relabeling, as required.

\paragraph{The full existential-second-order query statement.}
The remaining assertion is \(C(SF)=QNPTIME_{\rm tot}\) under the
explicit encoding, including arbitrary fixed answer rank and empty
universes. The forward direction guesses the fixed-arity auxiliary
relation table and evaluates the fixed first-order matrix. The guess has
polynomial size in the explicit universe size and formula evaluation has
a fixed polynomial exponent. It therefore recognizes the query graph in
nondeterministic polynomial time.

For the reverse direction fix a total consistent query in that class,
of answer rank \(b\). First work on nonempty universes. Expand the input
schema by a unary predicate \(U\), required to be the whole universe,
and unary predicates \(V_1,\ldots,V_b\), each required to be a singleton.
The expanded class consists of precisely those valid marked databases
whose marked tuple belongs to the query answer. It is closed under
isomorphism. Its canonical table encoding is recognized in
nondeterministic polynomial time: check the marker conditions, extract
their unique tuple, and run the original query recognizer on that
canonical database and tuple. The full unary \(U\) table makes this
an instance of the nonempty-input-type clause of
\cref{thm:p09-fagin74-theorem6}, even if the original schema is empty.
The explicit encoding conversions have polynomial overhead; no
binary-size-only encoding of a pure universe is substituted here.

That theorem gives an existential-second-order sentence defining the
expanded class. In its first-order matrix replace \(U(x)\) by
\(x=x\), and replace \(V_j(x)\) by \(x=d_j\), after renaming bound
variables to avoid the free output tuple \(\bar d\). This yields an
existential-second-order definition of exactly the required tuple query
on nonempty inputs. The function-elimination part of the cited theorem's
proof already converts auxiliary functions to relation graphs with
totality and uniqueness conditions, so only finitely many auxiliary
relations \(R_1,\ldots,R_l\) need be quantified. If there are none,
the resulting expression is already in \(F\subseteq SF\).

The source grammar permits one outer existential relation quantifier;
the following coding reduces this finite list to one. On a universe with
at least two elements choose distinct individual variables \(a,b\).
Use fixed, distinct tag words of a common length \(s\) over \(a,b\),
one per auxiliary relation, and let \(r\) be their maximum arity.
Replace
\[
 R_i(x_1,\ldots,x_{r_i})
 \quad\hbox{by}\quad
 R(\operatorname{tag}_i(a,b),x_1,\ldots,x_{r_i},a,\ldots,a),
\]
padding the final block to length \(r\). The displayed tag is just a
fixed tuple of the variables \(a,b\), so it is not an extra term-forming
operation. Distinct tags make the slices disjoint. Every family of
auxiliary relations extends to a single such \(R\) by filling those
slices, and every \(R\) determines a family by restriction. Thus
\(\exists R\exists a\exists b(a\ne b\land\theta_R)\) has exactly
the required meaning for these inputs, with the entire individual
quantification inside the first-order matrix of one \(SF\) expression.

There are only finitely many singleton diagrams for the fixed original
schema. Their correct answers form a first-order disjunction \(\chi\)
of complete diagrams; the unique candidate tuple needs no additional
choice. For an empty universe, positive answer rank forces the empty
answer, while rank zero has one fixed truth value \(\epsilon\). A
single first-order matrix under \(\exists R\) handles all cases by
disjoining: the accepted singleton diagrams, the at-least-two-elements
coding just constructed, and the empty-domain condition conjoined with
\(\epsilon\) when the answer rank is zero. The unused auxiliary
relation can always be chosen, including as the empty relation. This
is still one \(SF\) expression. It proves the total-query assertion
without equating partial functions having different definedness domains.

\end{proof}
\end{annexitem}

\begin{annexitem}{annex:p9.n-c-004.2.empty}{thm:p09-vardi82-empty-boundary}{Empty input and projection}
\begin{proof}
Use the literal grammar in \cref{def:p09-vardi82-languages}.
If a Boolean expression existentially projects even one
individual variable, it is false on an empty universe. Consequently an
expression true there has no such projected variable and no free output
variable. Every fixed-point component is then nullary: a positive-arity
component would require individual argument slots, and those slots would
have to be projected to obtain a Boolean answer.

A nullary positive operator acts on the two truth values. Its least
fixed point is true precisely when its first iterate from false is true:
if that first iterate is false it is already fixed; if true, monotonicity
keeps it true. Replace its nullary relation variable by a first-order
falsity sentence in the first-order matrix. The result is a first-order
sentence with exactly that least-fixed-point truth value. This also covers
the more permissive reading with a finite conjunction of nullary
fixed-point expressions: replace each component and conjoin the results.
The original first-order branch of the grammar is already first-order.
Thus every Boolean \(YF^+\) expression true on the empty universe is
first-order definable on all universes.

For a concrete obstruction, consider even cardinality on finite linear
orders, with zero counted as even. It is a total consistent polynomial-time
query under explicit universe representation. It is not first-order
definable, as the following finite game argument shows. With \(r\)
moves left, maintain that corresponding gaps between already selected
points have equal sizes or both have at least \(2^r-1\) elements.
Initially this holds for any two orders of sizes at least \(2^r-1\).
Equal finite gaps are answered at the same offset. In two large gaps,
put \(h=2^{r-1}-1\). If the chosen point has fewer than \(h\) points
on one side, match that short side exactly in the other gap; both remaining
long sides then have at least \(h\) points. The two sides cannot both
be short, because the original gap has at least \(2h+1\) points.
If neither side is short, choose a point with at least \(h\) points
on each side in the other gap. This maintains the invariant for \(r-1\)
moves. Repeated choices are matched to their already selected partners.
At zero moves the selected points have the same equality and order
relations, so the correspondence is a partial isomorphism.

For completeness, the link to formulas is the elementary induction on
quantifier rank: atomic formulas agree under a partial isomorphism;
Boolean connectives preserve agreement; for an existential formula choose
its witness as the next move and use the winning response together with
the induction hypothesis. Universal formulas follow by negation. Hence
orders of size at least \(2^r-1\) satisfy the same sentences of
quantifier rank at most \(r\). Choosing one even and one odd size above
that threshold contradicts any proposed first-order parity sentence.
Even cardinality, true on the empty order, is therefore outside this
literal all-domain \(YF^+\) grammar.

This does not contradict the nonempty-domain capture construction. On
nonempty orders its projected tableau defines parity. The closed language
\(Y\) may disjoin the empty-domain truth case with that projection.
It does show why the source-cited collapse \(C(Y)=C(YF^+)\) cannot be
imported into the present empty-domain convention without a corresponding
syntax or semantic qualification.
\end{proof}
\end{annexitem}
